% ============================================================================
% LANGENTWURF LE-022: Online-Banking
% Profil: S (Senioren 65+) | Tier: 2 (Standard)
% ============================================================================

\documentclass[11pt,a4paper]{article}
\usepackage[utf8]{inputenc}
\usepackage[T1]{fontenc}
\usepackage[ngerman]{babel}
\usepackage{geometry}
\usepackage{fancyhdr}
\usepackage{lastpage}
\usepackage{xcolor}
\usepackage{tcolorbox}
\usepackage{enumitem}
\usepackage{longtable}
\usepackage{hyperref}
\usepackage{amssymb}
\usepackage{eurosym}

\geometry{left=2.5cm, right=2.5cm, top=2.5cm, bottom=2.5cm}
\definecolor{fach}{HTML}{b35c00}
\definecolor{gray}{HTML}{666666}
\definecolor{successgreen}{HTML}{2a7a4b}
\definecolor{warnred}{HTML}{c62828}

\tcbuselibrary{skins,breakable}
\newtcolorbox{scorebox}{ colback=white, colframe=fach, fonttitle=\bfseries, title=Didaktik-Score, breakable }
\newtcolorbox{warnbox}[1][]{ colback=warnred!5, colframe=warnred, fonttitle=\bfseries, title=#1, breakable }
\newtcolorbox{infobox}[1][]{ colback=fach!5, colframe=fach, fonttitle=\bfseries, title=#1, breakable }

\pagestyle{fancy}
\fancyhf{}
\fancyhead[L]{\small\textcolor{gray}{Digitale Grundlagen -- 022 Online-Banking}}
\fancyhead[R]{\small\textcolor{gray}{LE Tier 2 / Profil S}}
\fancyfoot[C]{\small\thepage{} / \pageref{LastPage}}
\renewcommand{\headrulewidth}{0.4pt}

\begin{document}

\begin{titlepage}
    \centering
    \vspace*{2cm}
    {\Large\textcolor{gray}{Digitale Grundlagen -- Senioren 65+}}\\[2cm]
    {\Huge\textcolor{fach}{\textbf{Langentwurf}}}\\[0.5cm]
    {\large Tier 2 | Profil S}\\[2cm]
    {\LARGE\textbf{022 -- Online-Banking sicher nutzen}}\\[0.5cm]
    {\large Modul F: Sicherheit \& Finanzen}\\[2cm]
    \begin{tabular}{rl}
        \textbf{Zeitrahmen:} & 60--75 Minuten \\
        \textbf{Vorkenntnisse:} & Grundbedienung, Sicherheit (010) \\
    \end{tabular}
    \vfill
    {\normalsize Stand: \today}
\end{titlepage}

\tableofcontents
\newpage

\section{Bedingungsanalyse}

\subsection{Ausgangssituation}
\begin{itemize}
    \item Große Angst vor Online-Banking (``unsicher'')
    \item Angst, Geld zu verlieren oder Fehler zu machen
    \item Oft noch Papier-Überweisungen
    \item Unklarheit: Was ist sicher, was nicht?
    \item Banking-App vs. Browser unbekannt
\end{itemize}

\subsection{Typische Probleme}
\begin{itemize}
    \item ``Ist das wirklich sicher?''
    \item ``Was ist TAN, PIN, Passwort?''
    \item ``Was passiert, wenn ich einen Fehler mache?''
    \item ``Kann jemand mein Geld stehlen?''
    \item ``Soll ich die App oder den Browser nutzen?''
\end{itemize}

\section{Sachanalyse}

\subsection{Kernkonzepte}
\begin{enumerate}
    \item \textbf{Banking-App:} Offizielle App der eigenen Bank (sicherer als Browser)
    \item \textbf{PIN:} Persönliche Identifikationsnummer (nur du kennst sie)
    \item \textbf{TAN:} Transaktionsnummer -- einmalig gültig für eine Überweisung
    \item \textbf{Zwei-Faktor:} Zwei verschiedene Wege zur Bestätigung
    \item \textbf{Phishing:} Betrugsversuche per E-Mail/SMS
\end{enumerate}

\subsection{Alltagsanalogien}
\begin{tabular}{|l|p{8cm}|}
\hline
\textbf{Begriff} & \textbf{Analogie} \\
\hline
PIN & Schlüssel zu deinem Schließfach \\
TAN & Einmalcode wie beim Paket abholen \\
Zwei-Faktor & Haustür + Tresor (zwei Schlösser) \\
Phishing & Betrüger am Telefon, der sich als Bank ausgibt \\
Banking-App & Bankfiliale in der Hosentasche \\
\hline
\end{tabular}

\section{Didaktische Analyse}

\subsection{Stellung in der Reihe}
Aufbauend auf Sicherheitsgrundlagen (010). Erste Einheit im Modul ``Sicherheit \& Finanzen''.

\subsection{Didaktische Reduktion}
\textbf{Ausgeklammert:} Wertpapierhandel, Kreditkarten, Auslandsüberweisungen, SEPA-Details.

\textbf{Fokus:} Sicherheitsgefühl stärken, einfache Überweisungen verstehen, Phishing erkennen.

\section{Lernziele}

\begin{tabular}{|c|p{7cm}|c|c|}
\hline
\textbf{Nr.} & \textbf{Die Teilnehmer können...} & \textbf{Stufe} & \textbf{Niveau} \\
\hline
FZ1 & die Banking-App ihrer Bank finden & Anwenden & Basis \\
FZ2 & den Unterschied PIN/TAN erklären & Verstehen & Basis \\
FZ3 & Phishing-Versuche erkennen (3 Merkmale) & Analysieren & Basis \\
FZ4 & den Kontostand prüfen & Anwenden & Basis \\
FZ5 & eine Überweisung nachvollziehen & Verstehen & Vertiefung \\
\hline
\end{tabular}

\section{Methodische Analyse}

\begin{warnbox}[5 Sicherheitsregeln -- NIEMALS vergessen!]
\begin{enumerate}
    \item \textbf{Niemals} PIN/TAN am Telefon nennen -- die Bank fragt nie danach!
    \item \textbf{Niemals} auf Links in E-Mails klicken, die zur Bank führen sollen
    \item \textbf{Immer} die App nutzen oder Adresse selbst eintippen
    \item \textbf{Immer} Kontostand regelmäßig prüfen
    \item \textbf{Sofort} die Bank anrufen bei Verdacht
\end{enumerate}
\end{warnbox}

\begin{infobox}[Banking-App vs. Browser]
\begin{tabular}{|l|l|l|}
\hline
& \textbf{Banking-App} & \textbf{Browser} \\
\hline
Sicherheit & Sehr hoch & Hoch (wenn richtige Seite) \\
Bedienung & Einfacher & Komplizierter \\
Empfehlung & \textbf{Ja} (bevorzugt) & Nur am Computer \\
\hline
\end{tabular}
\end{infobox}

\section{Verlaufsplanung}

\begin{longtable}{|p{1cm}|p{2cm}|p{5cm}|p{3cm}|p{2cm}|}
\hline
\textbf{Zeit} & \textbf{Phase} & \textbf{Inhalt} & \textbf{TN-Aktivität} & \textbf{Material} \\
\hline
0--10 & Einstieg & ``Was ist Ihre größte Sorge beim Online-Banking?'' Ängste ernst nehmen & Berichten & -- \\
\hline
10--25 & Grundlagen & PIN, TAN, Zwei-Faktor erklären mit Analogien & Zuhören, Fragen & S.1 \\
\hline
25--40 & Sicherheit & 5 Sicherheitsregeln, Phishing-Beispiele zeigen & Notieren & S.1 \\
\hline
40--55 & Banking-App & App der eigenen Bank finden, Kontostand prüfen & Demo beobachten & S.2 \\
\hline
55--70 & Überweisung & Aufbau einer Überweisung erklären (IBAN, Betrag, Verwendungszweck) & Nachvollziehen & S.2-3 \\
\hline
\end{longtable}

\section{Lösungen}

\textbf{PIN vs. TAN:}
\begin{itemize}
    \item \textbf{PIN:} Persönliche Geheimnummer, bleibt immer gleich, öffnet den Zugang
    \item \textbf{TAN:} Einmalcode, nur für eine Transaktion gültig, bestätigt Überweisung
\end{itemize}

\textbf{Phishing erkennen (3 Merkmale):}
\begin{itemize}
    \item \textbf{Dringlichkeit:} ``Sofort handeln, sonst wird Konto gesperrt!''
    \item \textbf{Links:} Aufforderung, auf einen Link zu klicken
    \item \textbf{Daten:} Frage nach PIN, TAN, Passwort
\end{itemize}

\textbf{Im Notfall:}
\begin{itemize}
    \item Sofort Bank anrufen (Nummer auf der Karte!)
    \item Karte sperren lassen: \textbf{116 116} (bundesweit, kostenlos)
    \item Keine weiteren Überweisungen tätigen
    \item Polizei informieren bei Betrug
\end{itemize}

\section{Didaktik-Score}

\begin{scorebox}
\begin{tabular}{|c|l|c|c|p{3.5cm}|}
\hline
\# & Dimension & Gew. & Score & Notiz \\
\hline
1 & Differenzierung & 15\% & 8/10 & Basis/Vertiefung ok \\
2 & Taxonomie & 10\% & 9/10 & Verstehen + Analysieren \\
3 & Alltagsrelevanz & 10\% & 10/10 & Sehr hohe Relevanz \\
4 & Aktivierung & 10\% & 8/10 & Viel Demo, wenig Praxis \\
5 & Scaffolding & 10\% & 9/10 & Klare Regeln \\
6 & Feedback & 10\% & 8/10 & Ermutigung wichtig \\
7 & Sprachliche Klarheit & 10\% & 9/10 & Einfach erklärt \\
8 & Motivation & 10\% & 10/10 & Angst abbauen = motivierend \\
9 & Barrierefreiheit & 10\% & 9/10 & Große Schrift \\
10 & Praktikabilität & 5\% & 8/10 & 70 Min gut \\
\hline
\multicolumn{3}{|r|}{\textbf{GESAMT:}} & \textbf{9.05/10} & \textcolor{successgreen}{\textbf{FREIGABE}} \\
\hline
\end{tabular}
\end{scorebox}

\end{document}
